% GENERATED FILE. Edit canonical lessons, then run npm run generate:books.
% canonical-source-hash: fnv1a64:d87577d80ba17ccb
% canonical-lessons: MR-C294-pur, MR-C294-sacne, MR-C294-madatkarya, MR-C294-svayamsevak, MR-C294-adakne

\chapter{The Flood}
\label{ch:mr-the-flood}
\hlchaptermodality{294}

\begin{chapteropening}
\textbf{By the end of this chapter:} \emph{I can describe a flood: rising water, people stuck, volunteers helping.}

Complete the last lesson of chapter 294: I can describe a flood: rising water, people stuck, volunteers helping.
\end{chapteropening}

\section[\texorpdfstring{\mr{पूर} (pūr)}{पूर (pūr)}]{\mr{पूर} (pūr) — a flood}

\label{lesson:MR-C294-pur}



\begin{quote}
Before the new one: say the Marathi for to cave in, then the Marathi for a landslide.
\end{quote}

\subsection*{You'll want to know: \mr{पूर}}
\textbf{\mr{पूर}} — \emph{pūr} — \textquotedblleft{}a flood\textquotedblright{}.

\begin{quote}
\mr{नदीला} \mr{पूर} \mr{आल्यास} \mr{हा} \mr{पूल} \mr{बंद} \mr{होईल}.
\end{quote}

\begin{quote}
\emph{nadīlā pūr ālyās hā pūl band hoīl.}
\end{quote}

\begin{quote}
Should the river flood, this bridge will be closed.
\end{quote}

Masculine.

\mr{नदीला} \mr{पूर} \mr{येणे}, the river floods; \mr{आल्यास}, should it come.

\begin{cousinweb}
\textbf{\mr{पूर}} (\emph{pūr}) goes back to Sanskrit \emph{pūra}, 'filling; a flood': CDIAL lists the Marathi word, as \emph{purā}, 'filling; flood', under it (CDIAL 8330).
\end{cousinweb}

\begin{grammarlens}[title={What you've built}]
One more word for telling what happened.
\end{grammarlens}

\subsection*{Your turn}
\begin{itemize}
\raggedright
  \item \emph{Say it:} \emph{pūr}
  \item \emph{Say it:} \emph{pūr}, once more
  \item \emph{Say it:} \emph{nadīlā pūr ālyās hā pūl band hoīl.} — Should the river flood, this bridge will be closed.
  \item \emph{Recall:} say the Marathi for to cave in, then the Marathi for a landslide, then say \emph{pūr} again
\end{itemize}

\begin{tcolorbox}[breakable,colback=teal!4,colframe=teal!35!black,title={Before you move on}]
What does \mr{पूर} mean? (\textquotedblleft{}A flood\textquotedblright{}.) Say it once more.
\end{tcolorbox}
\section[\texorpdfstring{\mr{साचणे} (sācṇe)}{साचणे (sācṇe)}]{\mr{साचणे} (sācṇe) — to collect; to pile up}

\label{lesson:MR-C294-sacne}



\begin{quote}
Before the new one: say the Marathi for a landslide, then the Marathi for a flood.
\end{quote}

\subsection*{You'll want to know: \mr{साचणे}}
\textbf{\mr{साचणे}} — \emph{sācṇe} — \textquotedblleft{}to collect; to pile up\textquotedblright{}.

\begin{quote}
\mr{रस्त्यावर} \mr{गुडघाभर} \mr{पाणी} \mr{साचलं}.
\end{quote}

\begin{quote}
\emph{rastyāvar guḍghābhar pāṇī sāclaṁ.}
\end{quote}

\begin{quote}
Knee-deep water collected on the road.
\end{quote}

Past (he …): \mr{तो} \mr{साचला}, \emph{to sāclā}.

Of water, dust or work: \mr{काम} \mr{साचलं}, the work piled up.

\begin{cousinweb}
\textbf{\mr{साचणे}} (\emph{sācṇe}) is inherited: CDIAL lists Marathi \emph{sācṇe}, 'to accumulate', under \emph{*saṁcayayati}, 'piles up' (CDIAL 12867).
\end{cousinweb}

\begin{grammarlens}[title={What you've built}]
One more word for telling what happened.
\end{grammarlens}

\subsection*{Your turn}
\begin{itemize}
\raggedright
  \item \emph{Say it:} \emph{sācṇe}
  \item \emph{Say it:} \emph{sācṇe}, once more
  \item \emph{Say it:} \emph{rastyāvar guḍghābhar pāṇī sāclaṁ.} — Knee-deep water collected on the road.
  \item \emph{Recall:} say the Marathi for a landslide, then the Marathi for a flood, then say \emph{sācṇe} again
\end{itemize}

\begin{tcolorbox}[breakable,colback=teal!4,colframe=teal!35!black,title={Before you move on}]
What does \mr{साचणे} mean? (\textquotedblleft{}To collect; to pile up\textquotedblright{}.) Say it once more.
\end{tcolorbox}
\section[\texorpdfstring{\mr{मदतकार्य} (madatkārya)}{मदतकार्य (madatkārya)}]{\mr{मदतकार्य} (madatkārya) — relief work}

\label{lesson:MR-C294-madatkarya}



\begin{quote}
Before the new one: say the Marathi for a flood, then the Marathi for to collect.
\end{quote}

\subsection*{You'll want to know: \mr{मदतकार्य}}
\textbf{\mr{मदतकार्य}} — \emph{madatkārya} — \textquotedblleft{}relief work\textquotedblright{}.

\begin{quote}
\mr{पूर} \mr{ओसरल्यावर} \mr{मदतकार्य} \mr{सुरू} \mr{झालं}.
\end{quote}

\begin{quote}
\emph{pūr osralyāvar madatkārya surū jhālaṁ.}
\end{quote}

\begin{quote}
Once the flood went down, relief work began.
\end{quote}

Neuter.

Built from \mr{मदत}, help, and Sanskrit kārya, work.

\begin{cousinweb}
\textbf{\mr{मदतकार्य}} (\emph{madatkārya}) is a Marathi compound built from \emph{madat}, 'help', a loan usually traced to Arabic \emph{madad}, and Sanskrit \emph{kārya}, 'work, a task' (Monier-Williams).
\end{cousinweb}

\begin{grammarlens}[title={What you've built}]
One more word for telling what happened.
\end{grammarlens}

\subsection*{Your turn}
\begin{itemize}
\raggedright
  \item \emph{Say it:} \emph{madatkārya}
  \item \emph{Say it:} \emph{madatkārya}, once more
  \item \emph{Say it:} \emph{pūr osralyāvar madatkārya surū jhālaṁ.} — Once the flood went down, relief work began.
  \item \emph{Recall:} say the Marathi for a flood, then the Marathi for to collect, then say \emph{madatkārya} again
\end{itemize}

\begin{tcolorbox}[breakable,colback=teal!4,colframe=teal!35!black,title={Before you move on}]
What does \mr{मदतकार्य} mean? (\textquotedblleft{}Relief work\textquotedblright{}.) Say it once more.
\end{tcolorbox}
\section[\texorpdfstring{\mr{स्वयंसेवक} (svayaṁsevak)}{स्वयंसेवक (svayaṁsevak)}]{\mr{स्वयंसेवक} (svayaṁsevak) — a volunteer}

\label{lesson:MR-C294-svayamsevak}



\begin{quote}
Before the new one: say the Marathi for to collect, then the Marathi for relief work.
\end{quote}

\subsection*{You'll want to know: \mr{स्वयंसेवक}}
\textbf{\mr{स्वयंसेवक}} — \emph{svayaṁsevak} — \textquotedblleft{}a volunteer\textquotedblright{}.

\begin{quote}
\mr{शंभर} \mr{स्वयंसेवकांनी} \mr{लोकांना} \mr{जेवण} \mr{वाटलं}.
\end{quote}

\begin{quote}
\emph{śambhar svayaṁsevakānnī lokānnā jevaṇ vāṭlaṁ.}
\end{quote}

\begin{quote}
A hundred volunteers handed out food to people.
\end{quote}

Masculine.

A woman volunteer is \mr{स्वयंसेविका}.

\begin{cousinweb}
\textbf{\mr{स्वयंसेवक}} (\emph{svayaṁsevak}) is a modern compound built from Sanskrit \emph{svayam}, 'of oneself', and \emph{sevaka}, 'one who serves' (Monier-Williams): one who serves of his own accord.
\end{cousinweb}

\begin{grammarlens}[title={What you've built}]
One more word for telling what happened.
\end{grammarlens}

\subsection*{Your turn}
\begin{itemize}
\raggedright
  \item \emph{Say it:} \emph{svayaṁsevak}
  \item \emph{Say it:} \emph{svayaṁsevak}, once more
  \item \emph{Say it:} \emph{śambhar svayaṁsevakānnī lokānnā jevaṇ vāṭlaṁ.} — A hundred volunteers handed out food to people.
  \item \emph{Recall:} say the Marathi for to collect, then the Marathi for relief work, then say \emph{svayaṁsevak} again
\end{itemize}

\begin{tcolorbox}[breakable,colback=teal!4,colframe=teal!35!black,title={Before you move on}]
What does \mr{स्वयंसेवक} mean? (\textquotedblleft{}A volunteer\textquotedblright{}.) Say it once more.
\end{tcolorbox}
\section[\texorpdfstring{\mr{अडकणे} (aḍakṇe)}{अडकणे (aḍakṇe)}]{\mr{अडकणे} (aḍakṇe) — to get stuck; to be caught}

\label{lesson:MR-C294-adakne}



\begin{quote}
Before the new one: say the Marathi for relief work, then the Marathi for a volunteer.
\end{quote}

\subsection*{You'll want to know: \mr{अडकणे}}
\textbf{\mr{अडकणे}} — \emph{aḍakṇe} — \textquotedblleft{}to get stuck; to be caught\textquotedblright{}.

\begin{quote}
\mr{पाण्यात} \mr{दोन} \mr{मुलं} \mr{अडकली} \mr{होती}.
\end{quote}

\begin{quote}
\emph{pāṇyāt don mulaṁ aḍaklī hotī.}
\end{quote}

\begin{quote}
Two children were stuck in the water.
\end{quote}

Past (he …): \mr{तो} \mr{अडकला}, \emph{to aḍaklā}.

The place takes -\mr{त}: \mr{गर्दीत} \mr{अडकणे}, to be stuck in a crowd.

\begin{cousinweb}
\textbf{\mr{अडकणे}} (\emph{aḍakṇe}) is inherited: CDIAL lists Marathi \emph{aḍakṇe}, 'to stick fast', under \emph{*aḍ}, 'obstruct, stop' (CDIAL 187).
\end{cousinweb}

\begin{grammarlens}[title={What you've built}]
One more word for telling what happened.
\end{grammarlens}

\subsection*{Your turn}
\begin{itemize}
\raggedright
  \item \emph{Say it:} \emph{aḍakṇe}
  \item \emph{Say it:} \emph{aḍakṇe}, once more
  \item \emph{Say it:} \emph{pāṇyāt don mulaṁ aḍaklī hotī.} — Two children were stuck in the water.
  \item \emph{Recall:} say the Marathi for relief work, then the Marathi for a volunteer, then say \emph{aḍakṇe} again
\end{itemize}

\begin{tcolorbox}[breakable,colback=teal!4,colframe=teal!35!black,title={Before you move on}]
What does \mr{अडकणे} mean? (\textquotedblleft{}To get stuck; to be caught\textquotedblright{}.) Say it once more.
\end{tcolorbox}
